\section{Conclusion}
\label{sec:conclusion}

We introduce the \ara{} protocol and its surrounding ecosystem as a foundation for agent-native scientific communication.
Together, they address two structural failures of the PDF format: knowledge that narrative conventions discard (failed attempts, implicit configurations, unexplored branches) and specifications too underspecified to execute.
\ara{} resolves both by restructuring a research contribution as a machine-actionable artifact, one that is navigable, complete, and verifiable without human interpretation.

The broader motivation is a shift already underway: AI agents are becoming first-class participants in research workflows, not tools that assist humans but autonomous contributors that read, reproduce, and extend scientific work.
That transition demands infrastructure built around agents from the start.
\ara{} is the core abstraction of that ecosystem, a common substrate through which human and machine researchers alike publish, verify, and build on scientific knowledge.
